\documentclass[aspectratio=169,usenames,dvipsnames]{beamer}
\beamertemplatenavigationsymbolsempty
% \documentclass{beamer}

% UNCOMMENT FOR HANDOUTS
% \usepackage{pgfpages}
% \pgfpagesuselayout{2 on 1}[landscape,letterpaper,border shrink=5mm]
% \setbeameroption{show notes}

\mode<presentation>
\usetheme{Boadilla}
\usecolortheme{spruce}
\usepackage{../LizStyle,ulem}
\usepackage{color}
\usepackage[color,matrix,arrow]{xy}
% \input xy
% \xyoption{all}
\usepackage{tikz}
\usepackage{tikz-cd}
\tikzset{commutative diagrams/diagrams={ampersand replacement=\&}}

\AtBeginSection{\frame{\sectionpage}}


% ==========Command for putting comments and removing them=======
% Visible:
\newcommand{\InstructorNotes}[1]{\textit{\color{orange}#1}}
\newcommand{\InstructorList}[1]{
\begin{itemize}
	\color{orange}
	#1
\end{itemize}
}
% Removed:
% \renewcommand{\InstructorNotes}[1]{}
% \renewcommand{\InstructorList}[1]{}


\newcommand{\bvfill}{\vskip0pt plus 1filll}

\usepackage{multimedia}
\usepackage{graphbox}

% \usepackage{algpseudocode}



\graphicspath{ {Figures/} {../Figures/}{../../Figures/} }



\begin{document}
%\pdfoptionpdfminorversion 5
%\logo{\includegraphics[height=1.5cm]{dukeshield.pdf}}
\title[]{Ch 8.1: Decision Trees}
\subtitle{Lecture 24 - CMSE 381}
\date{Mon, Nov 4, 2024}



\institute[MSU-CMSE]{Michigan State University\\
			::\\
          Dept of Computational Mathematics, Science \& Engineering}
\author[Dr.~Munch]{Prof.~Elizabeth Munch}

\begin{frame}
 \titlepage


\end{frame}
%--------------------------
\begin{frame}[c]{Announcements}

\begin{columns}
\column{.45\textwidth}
\centering

	\textbf{Last time:}
	\begin{itemize}
		\item Cubic Splines
	\end{itemize}

	\vspace{.2in}

	\textbf{This lecture:}
	\begin{itemize}
	\item 8.1 Decision Trees
\end{itemize}

\textbf{Announcements:}
	\begin{itemize}
	\item HW \#7 Sun, 11/10 
	\item Projects

	\end{itemize}


\column{.45\textwidth}
\centering
\includegraphics[width = .8\textwidth]{Day01-Schedule-Test3}
\end{columns}

\end{frame}
%--- Next Frame ---%



% %=================
% \section{Last time}
% %=================



% \begin{frame}[c]{The idea with polynomial functions}

% Learn the model:
% \begin{equation*}
% 	y_i = \beta_0 + \beta_1 x_i + \beta_2 x_i^2 + \cdots + \beta_K x_i^K + \e_i
% \end{equation*}

% \begin{columns}
% \column{.25\textwidth}
% \centering
% My data:

% \begin{equation*}
% 	\begin{pmatrix}
% 		1 \\
% 		3 \\
% 		0 \\
% 		2 \\
% 		5 \\
% 		1 \\
% 		8
% 	\end{pmatrix}
% \end{equation*}
% \column{.45\textwidth}
% \centering
% What I pass to \texttt{LinearRegression} or \texttt{LogisticRegression} as features:
% \begin{equation*}
% 		\bordermatrix{ & \texttt{x} & \texttt{x}^2 & \texttt{x}^3 \cr
% 			 &\qquad  1 \qquad & \qquad  & \qquad \qquad\cr
% 			 &\qquad  3 \qquad &  &  \cr
% 			 &\qquad  0 \qquad &  &  \cr
% 			 &\qquad  2 \qquad &  &  \cr
% 			 &\qquad  5 \qquad &  &  \cr
% 			 &\qquad  1 \qquad &  &  \cr
% 			 &\qquad  8 \qquad &  &  \cr
% 			 } \qquad
% \end{equation*}

% \end{columns}

% \end{frame}
% %--- Next Frame ---%

% \begin{frame}[c]{The idea with step functions }

% 	Knots $z_1 \leq z_2 \leq \cdots \leq z_K$. Let $c_j(x) = \mathrm{I}(z_j \leq x \leq x_{j+1})$.
% 	Learn the model:
% 	\begin{equation*}
% 		y_i = \beta_0 + \beta_1 c_1(x_i) + \beta_2 c_2(x_i) + \cdots + \beta_K c_K(x_i) + \e_i
% 	\end{equation*}

% 	\begin{columns}
% 	\column{.25\textwidth}
% 	\centering
% 	My data:

% 	\begin{equation*}
% 		\begin{pmatrix}
% 			1 \\
% 			3 \\
% 			0 \\
% 			2 \\
% 			5 \\
% 			1 \\
% 			8
% 		\end{pmatrix}
% 	\end{equation*}
% 	\column{.45\textwidth}
% 	\centering
% 	What I pass to \texttt{LinearRegression} or \texttt{LogisticRegression} as features:
% 	\begin{equation*}
% 			\bordermatrix{ & \texttt{[-1,3)} & \texttt{[3,7)} & \texttt{[7,11)} \cr
% 				1 &\qquad  & \qquad  & \qquad \qquad\cr
% 				3 &\qquad   &  &  \cr
% 				0 &\qquad  &  &  \cr
% 				2 &\qquad   &  &  \cr
% 				5 &\qquad   &  &  \cr
% 				1 &\qquad   &  &  \cr
% 				8 &\qquad   &  &  \cr
% 				 } \qquad
% 	\end{equation*}

% 	\end{columns}
% \end{frame}
% %--- Next Frame ---%

% \begin{frame}[t]{Basis functions}


% Pick a collection of functions $b_i$ and learn the model:
% \begin{equation*}
% 	y_i = \beta_0 + \beta_1b_1(x_i) + \beta_2 b_2(x_i) + \cdots + \beta_K b_K(x_i) + \e_i
% \end{equation*}

% \begin{columns}
% \column{.25\textwidth}
% \centering
% My data:

% \begin{equation*}
% 	\begin{pmatrix}
% 		1 \\
% 		3 \\
% 		0 \\
% 		2 \\
% 		5 \\
% 		1 \\
% 		8
% 	\end{pmatrix}
% \end{equation*}
% \column{.45\textwidth}
% \centering
% What I pass to \texttt{LinearRegression} or \texttt{LogisticRegression} as features:
% \begin{equation*}
% 	\begin{pmatrix}
% 	  b_1(1) & b_2(1) & \cdots & b_K(1) \\
% 	  b_1(3) & b_2(3) & \cdots & b_K(3) \\
% 	  b_1(0) & b_2(0) & \cdots & b_K(0) \\
% 	  b_1(2) & b_2(2) & \cdots & b_K(2) \\
% 	  b_1(5) & b_2(5) & \cdots & b_K(5) \\
% 	  b_1(1) & b_2(1) & \cdots & b_K(1) \\
% 	  b_1(8) & b_2(8) & \cdots & b_K(8)
% 	\end{pmatrix}
% \end{equation*}

% \end{columns}

% \end{frame}
% %--- Next Frame ---%


%==========
\section{Decision Trees}
%==========

\begin{frame}[t]{Big idea}
\begin{columns}
\column{.45\textwidth}
\centering

\InstructorList{
  \item Previously, we've always had some form of function for prediction: $f(X_1,\cdots, X_p) = \hat Y$
  \item Now, we make a sequence of decisions to make a prediction, either regression or classification version 
}

\column{.45\textwidth}
\centering
\includegraphics[align = c, width = \textwidth]{binarytree-faces-w800}

\tiny 
Image: \url{https://marekbennett.com/2014/02/14/decision-tree/}
\end{columns}
\end{frame}
%--- Next Frame ---%

\begin{frame}[c]{Subset of Hitters data}

\begin{columns}
\column{.35\textwidth}
\centering
\includegraphics[width = \textwidth, align = c]{Hitters-TreeSlides.png}
\column{.45\textwidth}
\centering
\InstructorList{
  \item Remove observations missing salary values
	\item log transform salary for something closer to bell shape
	\item Goal: predict log salary (can reverse by returning $\exp(x)$ if model returns $x$)
}

\end{columns}

\end{frame}
%--- Next Frame ---%

\begin{frame}[c]{First decision tree example}

\begin{columns}
\column{.25\textwidth}
\centering
\includegraphics[width = \textwidth, align = c]{Hitters-TreeSlides-OnlyLog.png}
\column{.35\textwidth}
\footnotesize
\InstructorList{
  \item Top split assigns observations with $Years<4.5$ to left branch
	\item Return mean response for players with that.
	\item Predictions:
	\InstructorList{
	  \item  mean log salary is 5.107, so returns $\exp(5.107) = \$165.174 $ thousand dollars
		\item $5.999 \Rightarrow \$402,834$
		\item $6.740 \Rightarrow \$845,346$
	}
}

\column{.35\textwidth}
\centering
\includegraphics[width = \textwidth, align = c]{Fig8-1}
\end{columns}

\end{frame}
%--- Next Frame ---%





\begin{frame}[c]{Interpretation of example}

\begin{columns}
\column{.35\textwidth}
\centering

\includegraphics[width = \textwidth, align = c]{Fig8-1}



\column{.45\textwidth}
\centering
\InstructorList{
  \item Years most important factor for determining salary
	\item Players with less experience earn lower salaries than more experienced
	\item For the less experienced players, number of hits plays little role in salary
	\item For more experienced players, number of hits affects it
	\item Likely an oversimplification of real relationship, but easier to interpret and has nice graphical representation
}

\end{columns}

\end{frame}
%--- Next Frame ---%
{
\setbeamercolor{background canvas}{bg=Apricot!25}
\begin{frame}[t]{Coding a regression decision tree}
\begin{columns}
\column{.45\textwidth}
\centering
\column{.45\textwidth}
\centering
\end{columns}
\end{frame}
}
%--- Next Frame ---%

\begin{frame}[c]{Regions defined by the tree}

\begin{columns}
\column{.35\textwidth}
\centering
\includegraphics[width = \textwidth, align = c]{Fig8-1}
\column{.45\textwidth}
\centering
\InstructorList{
\item Point out tree is "upside down"
\item Leaves of the tree:
\InstructorList{
  \item $R_1 = \{X \mid Years<4.5\}$
	\item $R_2 = \{X \mid Years >= 4.5, Hits <117.5\}$
	\item $R_3 = \{X \mid Years >= 4.5, Hits >=117.5\}$
}
\item Other splits are called Internal Nodes
\item Segments that connect nodes called Branches or Edges
}

\end{columns}

\end{frame}
%--- Next Frame ---%
\begin{frame}[c]{Viewing Regions Defined by Tree}

	\begin{columns}
	\column{.35\textwidth}
	\centering
	\includegraphics[width = \textwidth, align = c]{Fig8-1}
	\column{.45\textwidth}
	\centering
	
	\includegraphics[width = \textwidth, align = c]{Fig8-2}
	\end{columns}
	
	\end{frame}
	%--- Next Frame ---%
\begin{frame}[c]{How do we actually get the tree? Two steps}

\begin{columns}
\column{.45\textwidth}
\centering
\begin{enumerate}
	\item We divide the predictor space -- that is, the set of possible values for $X_1 , X_2 , \cdots , X _p$ --- into $J$ distinct and non-overlapping regions, $R_1 , R_2 , \cdots, R_J $.
	\item For every observation that falls into the region $R_j$, we make the same prediction $=$ the mean of the response values for the training observations in $R_j $.
\end{enumerate}
\column{.45\textwidth}
\centering
\includegraphics[width = \textwidth, align = c]{Fig8-3b}
\end{columns}

\end{frame}
%--- Next Frame ---%



\begin{frame}[c]{Step 1: How do we decide on $R_j$s?}

\begin{columns}
\column{.45\textwidth}
\centering
\textbf{Goal:}

Find boxes $R_1,\cdots,R_J$ that minimize

\begin{equation*}
	\sum_{j=1}^J \sum_{i \in R_j} (y_i-\hat y_{R_j})^2
\end{equation*}
$\hat y_{R_j} = $ mean response for training observations in $j$th box

\column{.45\textwidth}
\centering
\InstructorList{
  \item Can't actually check every possible partition
	\item Instead, go for top-down greedy approach called \textit{Recursive binary splitting}
	\item Begins at top of tree with all data points
	\item Uses best split at every step
}

\end{columns}

\end{frame}
%--- Next Frame ---%
\begin{frame}[t]{Recursive Binary Splitting}

\begin{columns}
\column{.45\textwidth}
\centering
\begin{itemize}
	\item Pick $X_j$
  \item Pick $s$ so that splitting into $\{X \mid X_j < s\}$ and $\{X \mid X_j \geq s\}$ results in largest possible reduction in RSS
\end{itemize}
\column{.45\textwidth}
\centering
\begin{align*}
	R_1(j,s) = \{X \mid X_j < s\}\\
	R_2(j,s) = \{X \mid X_j \geq s\}\\
\end{align*}
\begin{equation*}
	\sum_{i \mid x_i \in R_1(j,s)}(y_i - \hat y _{R_1})^2
	+
	\sum_{i \mid x_i \in R_2(j,s)}(y_i - \hat y _{R_2})^2
\end{equation*}
\end{columns}

\end{frame}
%--- Next Frame ---%

\begin{frame}[c]{Rinse and repeat}

\begin{columns}
\column{.25\textwidth}
\centering
\InstructorList{
  \item Do this repeatedly
	\item Each time can split one of the previously identified regions
	\item Keep going until some stopping criterion is reached. E.g. until each region has at most 5 observations
}

\column{.35\textwidth}
\centering
\includegraphics[width = \textwidth, align = c]{Fig8-3c}
\column{.35\textwidth}
\centering
\includegraphics[width = \textwidth, align = c]{Fig8-3b}
\end{columns}

\end{frame}
%--- Next Frame ---%

\begin{frame}[c]{Pruning}

\begin{columns}
\column{.45\textwidth}
\centering
\includegraphics[width = \textwidth, align = c]{Fig8-4}
\column{.45\textwidth}
\centering
\InstructorList{
  \item Big trees leave you open to potential overfitting
	\item Could just stop building earlier, but that's short sighted
	\item Instead, grow a big tree and prune it back
	\item Find subtree with best test error rate
	\item Too many subtrees to test them all
}


\end{columns}

\end{frame}
%--- Next Frame ---%

% \begin{frame}[c]{Weakest Link Pruning }
% 	{Also called Cost complexity pruning}

% \begin{columns}
% \column{.45\textwidth}
% \centering

% For every $\alpha$, there is a subtree $T$ that minimizes:
% \begin{equation*}
% 	\sum_{m=1}^{|T|} \sum_{i \mid x_i \in R_m } (y_i - \hat y_{R_m})^2 + \alpha |T|
% \end{equation*}

% \begin{itemize}
% 	\item $|T| = $ number of terminal nodes of $T$
% 	\item $R_m$ is rectangle for $m$th terminal node
% 	\item $\hat y_{R_m}$ is mean of training observations in $R_m$
% \end{itemize}
% \column{.45\textwidth}
% \centering
% \InstructorList{
%   \item $\alpha=0$ gets entire tree
% 	\item Increasing $\alpha$ penalizes size of tree
% 	\item Branches pruned from tree in nested and predictable fashion
% 	\item Easy to get trees for all vales of $\alpha$
% 	\item Pick $\alpha $ via CV
% }

% \end{columns}

% \end{frame}
% %--- Next Frame ---%
% \begin{frame}[c]{Algorithm version}

% \begin{columns}
% \column{.65\textwidth}
% \centering
% \includegraphics[width = \textwidth, align = c]{Alg8-1}

% \end{columns}

% \end{frame}
% %--- Next Frame ---%

\begin{frame}[c]{Result of pruning}%{Messing with $\alpha$}

\begin{columns}
\column{.45\textwidth}
\centering
\includegraphics[width = \textwidth, align = c]{Fig8-5}
\column{.45\textwidth}
\centering
\includegraphics[width = .49\textwidth, align = c]{Fig8-4}
\includegraphics[width = .49\textwidth, align = c]{Fig8-1}
\end{columns}

\InstructorNotes{Result of pruning is the three leaf tree on the right. This example technically uses a more complex pruning method than the one described in the slides, but the general idea is the same.}
\end{frame}
%--- Next Frame ---%

%==========
\section{Classification Decision Tree}
%==========

\begin{frame}[c]{Basic idea}

\begin{columns}
\column{.35\textwidth}
\centering
\includegraphics[width = \textwidth, align = c]{Fig8-6c}

\column{.45\textwidth}
\centering
\begin{itemize}
	\item $\hat p_{mk} = $ proportion of training observations in $R_m$ from the $k$th class
	\item $E = 1-\max_{k}(\hat p_{mk})$
	\InstructorNotes{Fraction of training observations not in the most common class }
\end{itemize}
\InstructorList{
\item Example from heart data
\item Can't use RSS, need error rate
  \item Could decide splits by classification error rate
	\item Not necessarily the best, you can use other options but I'm going to skip that. 
}
\end{columns}

\end{frame}
%--- Next Frame ---%

% \begin{frame}[t]{Gini index}

% \begin{columns}
% \column{.45\textwidth}
% \centering
% \begin{equation*}
% 	G = \sum_{k=1}^K \hat p_{mk} (1-\hat p_{mk})
% \end{equation*}
% \column{.45\textwidth}
% \centering
% \InstructorList{
%   \item Measure of total variance across $K$ classes
% 	\item small value if all $\hat p _{mk}$'s close to zero or 1
% 	\item Small value means node contains mostly observations from one class
% }

% \end{columns}

% \end{frame}
% %--- Next Frame ---%
% \begin{frame}[t]{Entropy }

% \begin{columns}
% \column{.45\textwidth}
% \centering
% \begin{equation*}
% 	D = - \sum_{k=1}^K \hat p_{mk} \log \hat p _{mk}
% \end{equation*}
% \column{.45\textwidth}
% \centering
% \InstructorList{
%   \item Positive because $0 \leq \hat p_{mk} \leq 1$
% 	\item Near zero if $\hat p_{mk}$ all near 0 or near 1
% 	\item Small value if $m$th node has majority one class
% }

% \end{columns}

% \end{frame}
%--- Next Frame ---%

\begin{frame}[c]{Example}

\begin{columns}
\column{.75\textwidth}
\centering
\includegraphics[width = \textwidth, align = c]{Fig8-6a}

\end{columns}

\end{frame}
%--- Next Frame ---%
\begin{frame}[c]{Pruning the example }

\begin{columns}
\column{.45\textwidth}
\centering
\includegraphics[width = \textwidth, align = c]{Fig8-6b}
\column{.35\textwidth}
\centering
\includegraphics[width = \textwidth, align = c]{Fig8-6c}

\end{columns}

\end{frame}
%--- Next Frame ---%

{
\setbeamercolor{background canvas}{bg=Apricot!25}
\begin{frame}[t]{More coding!}
\begin{columns}
\column{.45\textwidth}
\centering
\column{.45\textwidth}
\centering
\end{columns}
\end{frame}
}
%--- Next Frame ---%
\begin{frame}[c]{Linear models vs trees }

\begin{columns}
\column{.45\textwidth}
\centering
\includegraphics[width = \textwidth, align = c]{Fig8-7-top}
\InstructorNotes{Obviously linear does better here}
\column{.45\textwidth}
\centering
\includegraphics[width = \textwidth, align = c]{Fig8-7-bottom}
\InstructorNotes{Not going to beat this case though}
\end{columns}

\end{frame}
%--- Next Frame ---%
\begin{frame}[t]{Pros/Cons}

\begin{columns}
\column{.45\textwidth}
\centering

\textbf{Pros:}

\InstructorList{
  \item Trees are very easy to explain to people. Often easier to explain than linear regression!
% \item Some people believe that decision trees more closely mirror human decision-making than do the regression and classification approaches seen in previous chapters.
\item Trees can be displayed graphically, and are easily interpreted even by a non-expert (especially if they are small).
\item  Trees can easily handle qualitative predictors without the need to create dummy variables.
}


\column{.45\textwidth}
\centering

\textbf{Cons:}

\InstructorList{
  \item Not as accurate as other methods of classification and regression
	\item Not robust: small change in data can cause large change in estimated tree
	\item Fix..... aggregate many decision trees
}


\end{columns}

\end{frame}
%--- Next Frame ---%

\begin{frame}[c]{TL:DR}

\begin{columns}
\column{.35\textwidth}
\centering
\begin{itemize}
	\item Split into regions by greedily decreasing RSS
	\item Prune tree by using cost complexity
	\item Not robust - Next time, figure out how to aggregate trees
\end{itemize}
\column{.3\textwidth}
\centering
\includegraphics[width = \textwidth, align = c]{Fig8-1}
\column{.3\textwidth}
\centering
\includegraphics[width = \textwidth, align = c]{Fig8-2}
\end{columns}

\end{frame}
%--- Next Frame ---%

\begin{frame}[t]{Next time}

\includegraphics[align = c, width = .5\textwidth]{Day01-Schedule-Test3}


\end{frame}
%--- Next Frame ---%


\end{document}
